\section{Tipos de Dados}

A forma comum em estatística de representar dados é por meio de matrizes retangulares
chamadas de \textbf{tabelas} -- ou também, devido aos softwares usados para tratá-las,
de \textbf{quadro de dados} (\textit{Dataframe}).
Por padrão, sendo $D$ um quadro de dados de dimensões $m \times n$, cada coluna de $D$
corresponde a um \textbf{atributo}, também chamado de variável ou \textit{feature}, e cada
linha corresponde a um registro, que representa um indivíduo na amostra coletada.

Cada atributo possui um domínio de valores possíveis, que embora possam ser os mais
variados possíveis, podem ser organizados numa taxonomia simples.
Os tipos de atributos podem ser divididos em dois grande grupos: \textbf{quantitativos} ou
\textbf{qualitativos}.
O primeiro refere-se a dados com domínio num conjunto numérico, sendo
dividido em

\begin{itemize}
\item \textbf{Discretos} -- São atributos que assumem valores em subconjuntos
dos números inteiros $\mathbb{Z}$, sendo que, na maioria dos contextos, são não negativos.
Geralemnte, representam contagens, como número de filhos e idades.

\item \textbf{Contínuos} -- O domínio do atributo é um intervalo dos números reais
$\mathbb{R}$.
Na prática, pelas limitações ontológicas de representar um número real, a realização desses
valores se limita a valores aproximados.
A diferença fundamental para os discretos é que o conjunto de valores possíveis é não
enumerável.
Exemplos incluem, principalmente, medidas, como preços, alturas e pesos.
\end{itemize}

Os dados qualitativos representam uma qualidade dentro de um conjunto finito de
possibilidades chamadas de fatores, sendo dividos em
\begin{itemize}
\item \textbf{Ordinais} -- os fatores possuem uma ordem natural.
Exemplos são graus de uma qualidade, como classe social (baixa, média e alta) ou
everidade de um exame médico.

\item \textbf{Nominais} -- os fatores não possuem ordem natural, sendo também chamados
de categóricos, pois representam categorias.
Exemplos são grupos, como sexo ou estado cívil.
Um tipo especial de variáveis nominais são as dicotômicas, cujos valores possíveis
são apenas dois e de caráter antagônico: sucesso ou fracasso, positivo ou negativo,
sim ou não, 0 ou 1, etc.
\end{itemize}

Alternativamente às classificações apresentadas, podemos tipificar atributos
através de escalas de medição de seus valores.
As escalas são similares às classificações apresentadas anteriormente, sendo elas:

\begin{itemize}
\item Escala Nominal -- para dados nessa escala, só podemos dizer que seus
valores são diferentes de outros, sendo usada para categorizar indivíduos.
Um exemplo simples é o sexo de uma pessoa: feminino ou masculino.
Essa escala não suporta operações matemáticas, e uma medida de centralidade
comum para resumir dados nessa escala é a moda.

\item Escala Ordinal -- nessa escala, os valores podem ser ordenados, e podemos
então dizer que, além de diferentes, um valor é maior do que o outro.
Essa escala tem sua estrutura preservada por operações que mantenham a ordem.
Um exemplo de variável ordinal é o nível de escolaridade: fundamental, médio
e superior. Medidas de tendência central comuns para dados nessa escala são a
moda e a mediana.

\item Escala Intervalar - essa escala possui uma origem arbitrária e necessita
de uma unidade de medida, sendo um exemplo a temperatura de um ambiente.
Podemos afirmar que valores são diferentes, maiores e quão o maior em relação
a outro, e transformações afim não alteram a estrutura dessa escala.
Podemos utilizar medidas como média, moda e mediana.

\item Escala de Razão -- nessa escala, existe uma origem absoluta, e podemos dizer quantas
vezes um valor maior do que outro através de razões.
Um exemplo de variável nessa escala é o peso de um indivíduo.
Transformações lineares preservam a estrutura dessa escala, e podemos utilizar
medidas como média, moda e mediana.
\end{itemize}

Em muitos casos, os atributos de um quadro de dados podem ser modelados enquanto variáveis
aleatórias.
É preciso, para tanto, que os atributos tenham representações numéricas válidas,
especialmente no caso dos atributos qualitativos, a exemplo dos ordinais, que podem
ser representados por uma enumeração de $1$ até $n$, onde $n$ indica o número de
fatores, e dos dicotômicos, que são podem ser simplesmente representados por $0$ e $1$.
A modelagem por variáveis aleatórias, visa, nas tarefas de análise exploratória,
explicar os atributos aplicando os resultados da teoria de distribuições de probabilidade.
Com essa abordagem, pode-se explicar aspectos como frequências (absolutas relativas ou
acumuladas), pontos de centralidade, dispersão, assimetria, valores atípicos, entre outros.
